\documentclass[10pt,a4paper]{amsart}
\usepackage{amsmath,amssymb,amsthm,graphicx}
\usepackage{verbatim}
\newenvironment{reference}{\comment}{\endcomment}
%\newenvironment{reference}{}{}
\newenvironment{slogan}{\comment}{\endcomment}
\newenvironment{history}{\comment}{\endcomment}

% For commutative diagrams we use Xy-pic
\usepackage[all]{xy}

% We use 2cell for 2-commutative diagrams.
\xyoption{2cell}
% Theorem environments.
%
\theoremstyle{plain}
\newtheorem{theorem}[subsection]{Theorem}
\newtheorem{proposition}[subsection]{Proposition}
\newtheorem{lemma}[subsection]{Lemma}

\theoremstyle{definition}
\newtheorem{definition}[subsection]{Definition}
\newtheorem{example}[subsection]{Example}
\newtheorem{exercise}[subsection]{Exercise}
\newtheorem{situation}[subsection]{Situation}

\theoremstyle{remark}
\newtheorem{remark}[subsection]{Remark}
\newtheorem{remarks}[subsection]{Remarks}

\numberwithin{equation}{subsection}

% Macros
%
\def\lim{\mathop{\mathrm{lim}}\nolimits}
\def\colim{\mathop{\mathrm{colim}}\nolimits}
\def\Spec{\mathop{\mathrm{Spec}}}
\def\Hom{\mathop{\mathrm{Hom}}\nolimits}
\def\Ext{\mathop{\mathrm{Ext}}\nolimits}
\def\SheafHom{\mathop{\mathcal{H}\!\mathit{om}}\nolimits}
\def\SheafExt{\mathop{\mathcal{E}\!\mathit{xt}}\nolimits}
\def\Sch{\mathit{Sch}}
\def\Mor{\mathop{\mathrm{Mor}}\nolimits}
\def\Ob{\mathop{\mathrm{Ob}}\nolimits}
\def\Sh{\mathop{\mathit{Sh}}\nolimits}
\def\NL{\mathop{N\!L}\nolimits}
\def\CH{\mathop{\mathrm{CH}}\nolimits}
\def\proetale{{pro\text{-}\acute{e}tale}}
\def\etale{{\acute{e}tale}}
\def\QCoh{\mathit{QCoh}}
\def\Ker{\mathop{\mathrm{Ker}}}
\def\Im{\mathop{\mathrm{Im}}}
\def\Coker{\mathop{\mathrm{Coker}}}
\def\Coim{\mathop{\mathrm{Coim}}}

% Boxtimes
%
\DeclareMathSymbol{\boxtimes}{\mathbin}{AMSa}{"02}

%
% Macros for moduli stacks/spaces
%
\def\QCohstack{\mathcal{QC}\!\mathit{oh}}
\def\Cohstack{\mathcal{C}\!\mathit{oh}}
\def\Spacesstack{\mathcal{S}\!\mathit{paces}}
\def\Quotfunctor{\mathrm{Quot}}
\def\Hilbfunctor{\mathrm{Hilb}}
\def\Curvesstack{\mathcal{C}\!\mathit{urves}}
\def\Polarizedstack{\mathcal{P}\!\mathit{olarized}}
\def\Complexesstack{\mathcal{C}\!\mathit{omplexes}}
% \Pic is the operator that assigns to X its picard group, usage \Pic(X)
% \Picardstack_{X/B} denotes the Picard stack of X over B
% \Picardfunctor_{X/B} denotes the Picard functor of X over B
\def\Pic{\mathop{\mathrm{Pic}}\nolimits}
\def\Picardstack{\mathcal{P}\!\mathit{ic}}
\def\Picardfunctor{\mathrm{Pic}}
\def\Deformationcategory{\mathcal{D}\!\mathit{ef}}
\title{A local ring with a unique prime ideal and a nonzero flat ideal}
\author{The Stacks Project Authors}
\date{stacks-project:examples.tex:local-snapshot-20261012; GNU FDL 1.2 or later}
\begin{document}
\maketitle
\noindent\textit{Editorial note.} All delivered mathematical text is the author's own, copied from the pinned Stacks Project chapter file examples.tex; no statement or proof step was rewritten, reordered or removed.\par\medskip
\noindent\textit{Editorial note.} The counted proof is the author's complete argument: The counted statement's missing final period is the author's own (verbatim). The literature attribution lives only in the skipped provenance sentence.\par\medskip
\noindent\textit{Editorial note.} The label lemma-zero-dimensional-flat-ideal names the counted statement, and the proof environment immediately following it in the source is delivered verbatim as its proof.\par\medskip
\noindent\textit{Editorial note.} The counted statement is carried by the author's own lemma environment in the chapter source.\par\medskip
\noindent\textit{Allowed standard prerequisites (editorial).}\begin{itemize}
\item Context span 2851-2861: R = k[X\_i, Y\_i]/(X\_i - Y\_i X\_\{i+1\}, Y\_i\textasciicircum{}2) with images x\_i, y\_i; the chain x\_i = y\_i x\_\{i+1\} = y\_i y\_\{i+1\} x\_\{i+2\} = ...; R has only one prime ideal, namely m = (x\_i, y\_i); the ideal I = (x\_i) is claimed flat
\item Context span 2863-2876: the annihilator of x\_i in R is the ideal (x\_1, x\_2, ..., y\_i, y\_\{i+1\}, ...); the module M generated by e\_i with relations e\_i = y\_i e\_\{i+1\} is flat as the colimit of copies of R with transition maps R --y\_1--> R --y\_2--> R --y\_3--> ... (colimits of flat modules with flat... the displayed transition system); the annihilator of e\_i in M equals the annihilator of x\_i in R; every element of M, respectively I, is f e\_i, respectively h x\_i, so the map M -> I sending e\_i to x\_i is an isomorphism and I is flat
\end{itemize}\medskip
nonzero flat ideal. Here is the construction. Let $k$ be a field.
Let $X_i, Y_i$, $i \geq 1$ be variables. Take
$R = k[X_i, Y_i]/(X_i - Y_i X_{i + 1}, Y_i^2)$. Denote $x_i$, resp.\ $y_i$
the image of $X_i$, resp.\ $Y_i$ in this ring. Note that
$$
x_i = y_i x_{i + 1} = y_i y_{i +1} x_{i + 2} =
y_i y_{i + 1} y_{i + 2} x_{i + 3} = \ldots
$$
in this ring. The ring $R$ has only
one prime ideal, namely $\mathfrak m = (x_i, y_i)$. We claim that
the ideal $I = (x_i)$ is flat as an $R$-module.

\medskip\noindent
Note that the annihilator of $x_i$ in $R$ is the ideal
$(x_1, x_2, x_3, \ldots, y_i, y_{i + 1}, y_{i + 2}, \ldots)$.
Consider the $R$-module $M$ generated by elements $e_i$, $i \geq 1$ and
relations $e_i = y_i e_{i + 1}$. Then $M$ is flat as it is the
colimit $\colim_i R$ of copies of $R$ with transition maps
$$
R \xrightarrow{y_1} R \xrightarrow{y_2} R \xrightarrow{y_3} \ldots
$$
Note that the annihilator of $e_i$ in $M$ is the ideal
$(x_1, x_2, x_3, \ldots, y_i, y_{i + 1}, y_{i + 2}, \ldots)$.
Since every element of $M$, resp.\ $I$ can be written as
$f e_i$, resp.\ $h x_i$ for some $f, h \in R$ we see that the
map $M \to I$, $e_i \to x_i$ is an isomorphism and $I$ is flat.

\begin{lemma}
\label{lemma-zero-dimensional-flat-ideal}
\begin{slogan}
Zero dimensional ring with flat ideal.
\end{slogan}
There exists a local ring $R$ with a unique prime ideal
and a nonzero ideal $I \subset R$ which is a flat $R$-module
\end{lemma}

\begin{proof}
See discussion above.
\end{proof}

\end{document}
